\section{Related work}
\label{sec:related}

Self-correction studies ask whether models can fix errors when told to look for them, and largely answer no. Models asked to review their own answers rarely improve them \citep{huang2024}, struggle to locate errors even when told one exists \citep{tyen2024}, and remain blind to injected errors in their own generations even in benchmarks built to elicit correction \citep{tsui2025}. Prompted-critique studies find the same leniency toward supplied reasoning: models grading step-by-step solutions accept a large majority of wrong steps \citep{reflect2026,vair2026}, including when a single step was surgically edited \citep{breakingchain2026}. Our question is upstream of all of these: not whether a model can find an error when asked, but whether it notices one it was never told about, in reasoning it believes is its own.

The nearest work to ours measures unprompted checking, and we differ from it on the variable that turns out to matter. The Self-Verification Dilemma \citep{svd2026} finds that models rarely re-verify steps in naturally produced traces, framed as a general token-budget tradeoff, on open models. We plant audited errors rather than relying on natural ones, which fixes the ground truth; we separate readable from computed content, which locates the failure precisely; and we test frontier systems, which shows the readable side improving while the computed side does not. A cross-conversation analog \citep{reclaim2026} finds models re-derive conclusions rather than trusting remembered ones across sessions; our result is the within-trace complement, where the same deference that is prudent across sessions becomes total within one.

A mechanistic line of work converges on our deference reading from inside the network. Probing studies find that models internally represent whether an arithmetic step is correct even while their behavior ignores it \citep{validationgap2025}, which matches what we observe behaviorally: the failure is not missing information but an unused check. Relatedly, models executing long programs degrade as traces lengthen \citep{execdecay2026}; we hold length fixed and vary only the computation a check requires, isolating verification cost from execution burden.

Trained verifiers are the deliberate foil for our result. Process reward models supervise individual reasoning steps and readily learn to catch computation errors \citep{lightman2023}, so the checking ability is plainly learnable. Our finding is that the generating model does not deploy it unprompted, at any scale we test, and that neither instruction nor demonstration turns it on at inference time. The gap between what a dedicated verifier catches and what the generator checks by default is exactly the surface our measurement exposes.

Finally, our elicitation borrows from work on forced continuation, where models are made to continue from imposed prefixes to study safety recovery \citep{qi2025,vonrecum2026}. We use the same mechanism for measurement rather than attack: continuation from the model's own perturbed trace is what lets us observe the unprompted default, and error compounding in multi-step reasoning \citep{dziri2023} is why that default matters downstream.
